\section{Truth and Reality}

\begin{quotex}
And the independent reality of the individual, when we examine it, is in truth mere illusion. Apart from the community, what are separate men? It is the common mind within him which gives reality to the human being, and taken by himself, whatever else he is, he is no human ... If this is true of the social consciousness in its various forms, it is true certainly no less of that common mind which is more than social. The finite minds that in and for religion form one spiritual whole have indeed in the end no visible embodiment, and yet, except as members in an invisible community, they are nothing real. For religion, in short, if the one indwelling spirit is removed, there are no spirits left.

\flright{\textsc{Francis Herbert Bradley}, \textit{Truth and Reality}}

\end{quotex}





\flrightit{Posted on 2012-07-30 by Cologero}

\begin{center}* * *\end{center}

\begin{footnotesize}
\begin{sffamily}
\texttt{zero is zero on 2012-07-31 at 21:54 said:}

``For Julius Evola, on the other hand, ultimate reality is the Absolute Individual.'' I think he would only have called it that during his early philosophical period which he later rejected in favor of `Tradition'. Evola's later position on the question of individuality is in chapter 16 of Ride The Tiger and seems very close to Guenon's in the letter you are referring to.


\hfill

\end{sffamily}
\end{footnotesize}

